% Unit 114 terjemahan Bahasa Indonesia dari chapter8.tex baris 1353--1482.
% Sumber: Methods of Algebra, Volume 2, commit
% 9a5803ff77dd3257484cb177f851a73770a59dd3 (CC BY 4.0).
% stable-unit-id: o014.aljabr2.chapter8.bar-construction-b
% Cakupan: Bagian 8.7 dari makna resolusi bar hingga contoh homologi grup.

% segment-id: o014.li2.chapter8.bar-construction-b.p001
Objek simplisial teraugmentasi dalam
\eqref{eqn:bar-simplicial-resolution} yang diberikan oleh konstruksi bar
dapat dipandang secara heuristik sebagai pengganti $M$. Untuk
memperjelas gagasan ini, pertimbangkan pertanyaan berikut. Misalkan
$\Gamma$ sebuah grup dan lengkapi himpunan satu titik $\{\mathrm{pt}\}$
dengan aksi kanan trivial $\Gamma$. Bagaimana kita harus memahami
% segment-id: o014.li2.chapter8.bar-construction-b.q001
\[ \text{``ruang hasil bagi''}\; \{\mathrm{pt}\} / \Gamma ? \]
Dengan definisi naif hasil bagi himpunan atau ruang topologis, jawabannya
tentu tidak menarik; secara umum, mengambil hasil bagi oleh suatu aksi
yang tidak bebas akan menghapus informasi dari ruang semula. Akan tetapi,
berbagai kebutuhan dalam praktik mendorong kita mencari definisi yang
lebih halus bagi hasil bagi oleh aksi yang tidak bebas.

% segment-id: o014.li2.chapter8.bar-construction-b.p002
Untuk itu, tuliskan $\cated{Set}\Gamma$ untuk kategori semua himpunan
kanan-$\Gamma$ kecil, dan tinjau pasangan adjoin bebas--lupa
$\cate{Set}\leftrightarrows\cated{Set}\Gamma$. Konstruksi bar menghasilkan
dari himpunan kanan-$\Gamma$ satu titik $\{\mathrm{pt}\}$ sebuah objek
simplisial dalam $\cated{Set}\Gamma$; objek ini tepat himpunan simplisial
$\mathrm{E}\Gamma$ yang diperkenalkan dalam
Contoh~\ref{eg:classifying-space}. Verifikasi terperincinya diserahkan
sebagai latihan bab ini. Setelah fakta tersebut dipahami dan konstruksi bar
dipandang sebagai pengganti $\{\mathrm{pt}\}$ dalam situasi ini,
masuk akal untuk mendefinisikan $\{\mathrm{pt}\}/\Gamma$ sebagai
$(\mathrm{E}\Gamma)/\Gamma\simeq\mathrm{B}\Gamma$. Ini merupakan hasil
bagi oleh aksi bebas dan memberikan ruang klasifikasi $\Gamma$;
hasilnya lebih wajar dan biasanya lebih berguna daripada hasil bagi
himpunan yang naif.

% segment-id: o014.li2.chapter8.bar-construction-b.p003
Dalam arti apa objek simplisial teraugmentasi dalam
Contoh~\ref{eg:bar-adjunction} dan~\ref{eg:bar-comonad} dapat dikatakan
memberikan semacam pengganti, atau ``resolusi'', bagi objek semula?
Gagasannya berasal dari topologi, dan makna tepatnya melibatkan sifat
dapat dikontraksi dalam Definisi~\ref{def:contractible-aug}. Salah satu
hasil yang berkaitan adalah sebagai berikut.

% segment-id: o014.li2.chapter8.bar-construction-b.pr001
\begin{proposition}\label{prop:bar-contractible}
	Dalam situasi Contoh~\ref{eg:bar-adjunction}, misalkan diberikan pasangan
	adjoin
	% segment-id: o014.li2.chapter8.bar-construction-b.q002
	\[\begin{tikzcd}
		F : \mathcal{C} \arrow[shift left, r] & \mathcal{D}: G \arrow[shift left, l] ,
	\end{tikzcd} \quad \begin{array}{ll}
		\eta: \identity_{\mathcal{C}} \to GF & \text{unit} \\
		\epsilon: FG \to \identity_{\mathcal{D}} & \text{kounit}
	\end{array}\]
	dan bentuk objek simplisial teraugmentasi $\mathrm{Bar}(F,G)$ dalam
	$\End(\mathcal{D})$. Maka $G\mathrm{Bar}(F,G)$ dapat dikontraksi dari
	kiri.
\end{proposition}
\begin{proof}
	Gunakan notasi dari Contoh~\ref{eg:bar-adjunction}, misalnya $T=GF$ dan
	$\mu:T^2\to T$. Ingat bahwa
	$G\mathrm{Bar}(F,G)_n=T^{n+1}G$. Untuk setiap $n\geq-1$, definisikan
	% segment-id: o014.li2.chapter8.bar-construction-b.q003
	\[ k_n := \eta T^{n+1}G: T^{n+1}G \to T^{n+2}G . \]
	Ingat pula bahwa morfisme muka
	$d_i:T^{n+1}G\to T^nG$ pada $G\mathrm{Bar}(F,G)$ ialah
	$T^i\mu T^{n-i-1}G$ jika $0\leq i<n$, dan $T^nG\epsilon$ jika
	$i=n$; morfisme augmentasinya ialah $G\epsilon:TG\to G$
	(perhatikan bahwa notasi ini berbeda dari notasi dalam
	Definisi~\ref{def:contractible-aug}). Syarat dapat dikontraksi dari kiri
	kini diverifikasi satu demi satu.

	Pertama, identitas segitiga bagi adjungsi memberikan
	$(G\epsilon)k_{-1}=(G\epsilon)(\eta G)=\identity_G$.

	Kedua, $(T,\mu,\eta)$ adalah monad, sehingga
	$\mu(\eta T)=\identity_T$. Jika $n\geq0$, komposisikan kedua ruas di
	sebelah kanan dengan $T^nG$ untuk memperoleh
	$d_0k_n=\mu T^nG\cdot\eta T^{n+1}G=\identity_{T^{n+1}G}$.

	Terakhir, berdasarkan naturalitas $\eta$, diagram berikut komutatif bila
	$1\leq i\leq n+1$:
	% segment-id: o014.li2.chapter8.bar-construction-b.q004
	\[\begin{tikzcd}[column sep=large]
		T^{n+1} G \arrow[r, "{k_n = \eta T^{n+1}G}" inner sep=0.6em] \arrow[d, "{d_{i-1} = \varphi}"'] & T^{n+2}G \arrow[d, "{d_i = T\varphi}"] \\
		T^n G \arrow[r, "{k_{n-1} = \eta T^n G}"' inner sep=0.6em] & T^{n+1} G
	\end{tikzcd} \quad
		\varphi := \begin{cases}
			T^{i-1} \mu T^{n-i} G, & 1 \leq i \leq n \\
			T^n G\epsilon, & i = n+1.
		\end{cases}\]
	Ini memberikan $d_1k_0=k_{-1}(G\epsilon)$ ketika $i=1$ dan $n=0$,
	serta $d_ik_n=k_{n-1}d_{i-1}$ ketika $n\geq1$. Pernyataan terbukti.
\end{proof}

% segment-id: o014.li2.chapter8.bar-construction-b.p004
Walaupun kontraktibilitas
$G\mathrm{Bar}(F,G)$ tidak selalu mengakibatkan
$\mathrm{Bar}(F,G)$ dapat dikontraksi,
Proposisi~\ref{prop:bar-contractible} sudah memadai untuk penerapan
selanjutnya. Latihan bab ini akan memperkenalkan hasil-hasil lain mengenai
sifat dapat dikontraksi.

% segment-id: o014.li2.chapter8.bar-construction-b.p005
Dalam kasus kategori aditif, makna ``resolusi'' juga dapat
diangkat ke tingkat kompleks rantai dan dengan demikian dipahami
melalui aljabar linear. Mula-mula, menurut
Definisi~\ref{def:augmented-C-cplx}, objek simplisial teraugmentasi $X$
dalam kategori aditif memberikan augmentasi $\mathrm{C}^{\mathrm{aug}}X$
dari kompleks rantai $\mathrm{C}X$. Jika morfisme identitas suatu kompleks
rantai homotopik terhadap nol, kompleks rantai itu sendiri disebut
	\emph{homotopik nol}.

% segment-id: o014.li2.chapter8.bar-construction-b.th001
\begin{theorem}\label{prop:bar-null-homotopic}
	Misalkan $\mathcal{A}$ kategori aditif dan diberikan pasangan adjoin
	% segment-id: o014.li2.chapter8.bar-construction-b.q005
	\[\begin{tikzcd}
		F : \mathcal{A} \arrow[shift left, r] & \mathcal{B}: G, \arrow[shift left, l]
	\end{tikzcd}\]
	dan bentuklah objek simplisial teraugmentasi
	$\mathrm{Bar}(F,G)$ dalam $\End(\mathcal{B})$ (atau
	$G\mathrm{Bar}(F,G)$ dalam $\mathcal{A}^{\mathcal{B}}$).
	\begin{itemize}
		\item Perhatikan bahwa $\mathcal{A}^{\mathcal{B}}$ kategori aditif
		(lihat bagian pertama \S\ref{sec:Abel-cat-def}). Kompleks rantai
		$\mathrm{C}^{\mathrm{aug}}\bigl(G\mathrm{Bar}(F,G)\bigr)$ dari
		Definisi~\ref{def:augmented-C-cplx} homotopik nol.
		\item Untuk setiap $Y\in\Obj(\mathcal{B})$, evaluasi suku demi suku
		$\mathrm{Bar}(F,G)$ (atau $G\mathrm{Bar}(F,G)$) pada $Y$. Ini menghasilkan objek simplisial teraugmentasi
		$\mathrm{Bar}(Y)$ dalam $\mathcal{B}$ (atau
		$G\mathrm{Bar}(Y)$ dalam $\mathcal{A}$), dan kompleks rantai
		$\mathrm{C}^{\mathrm{aug}}\bigl(G\mathrm{Bar}(Y)\bigr)$ dalam
		$\mathcal{A}$ homotopik nol.
	\end{itemize}
\end{theorem}
\begin{proof}
	Proposisi~\ref{prop:bar-contractible} menunjukkan bahwa
	$G\mathrm{Bar}(F,G)$ dapat dikontraksi dari kiri. Karena syarat dapat
	dikontraksi bersifat absolut, $G\mathrm{Bar}(Y)$ yang diperoleh melalui
	funktor evaluasi $\mathrm{ev}_Y:\mathcal{A}^{\mathcal{B}}\to\mathcal{A}$
	secara alami juga dapat dikontraksi dari kiri.
	Terapkan Proposisi~\ref{prop:contractible-null-homotopic} pada kedua
	objek tersebut.
\end{proof}

% segment-id: o014.li2.chapter8.bar-construction-b.p006
Sekarang misalkan $\mathcal{A}$ dan $\mathcal{B}$ keduanya kategori
aditif. Untuk pasangan adjoin seperti dalam
Teorema~\ref{prop:bar-null-homotopic} dan
$Y\in\Obj(\mathcal{B})$, karena $\mathrm{Bar}(Y)_{-1}=Y$, dari
$\mathrm{C}^{\mathrm{aug}}(\mathrm{Bar}(Y))$ diperoleh morfisme
dalam $\cate{Ch}_{\geq0}(\mathcal{B})$
% segment-id: o014.li2.chapter8.bar-construction-b.q006
\begin{equation}\label{eqn:bar-resolution}
	\epsilon_Y: \mathrm{C}(\mathrm{Bar}(Y)) \to Y;
\end{equation}
morfisme ini kanonik dalam $Y$, dan bagian berderajat nolnya tepat morfisme
augmentasi $\epsilon_Y:FG(Y)\to Y$ dari $\mathrm{Bar}(Y)$, sehingga
notasinya sesuai.

% segment-id: o014.li2.chapter8.bar-construction-b.co001
\begin{corollary}\label{prop:bar-exactness-G}
	Misalkan diberikan pasangan adjoin antara kategori abelian\footnote{Di sini $F$ dan
	$G$ tentu funktor aditif; lihat Korolari~\ref{prop:automatic-additivity}.}
	% segment-id: o014.li2.chapter8.bar-construction-b.q007
	$\begin{tikzcd}
		F : \mathcal{A} \arrow[shift left, r] & \mathcal{B}: G \arrow[shift left, l]
	\end{tikzcd}$
	dan $Y\in\Obj(\mathcal{B})$. Morfisme yang diperoleh dengan menerapkan
	$G$ suku demi suku pada \eqref{eqn:bar-resolution}, yaitu
	% segment-id: o014.li2.chapter8.bar-construction-b.q008
	\[ G\epsilon_Y: \mathrm{C}(G\mathrm{Bar}(Y)) \to GY ,\]
	merupakan kuasi-isomorfisme dalam $\cate{Ch}(\mathcal{A})$.
\end{corollary}
\begin{proof}
	Teorema~\ref{prop:bar-null-homotopic} mengakibatkan
	$\mathrm{C}^{\mathrm{aug}}(G\mathrm{Bar}(Y))$ asiklik, dan hasilnya
	langsung mengikuti.
\end{proof}

% segment-id: o014.li2.chapter8.bar-construction-b.eg001
\begin{example}[Resolusi bar modul serta homologi/kohomologi Hochschild]\label{eg:HH-comonad}
	\index{resolusi bar@resolusi bar}
	\index[sym1]{BR}
	\index[sym1]{B'R}
	\index[sym1]{HH}
	Misalkan $\Bbbk$ gelanggang komutatif dan $R$ sebuah aljabar-$\Bbbk$.
	Tuliskan $\otimes:=\otimes_{\Bbbk}$. Menurut
	Contoh~\ref{eg:comonad-change-ring}, komonad yang ditentukan oleh
	pasangan adjoin
	% segment-id: o014.li2.chapter8.bar-construction-b.q009
	\[\begin{tikzcd}
		R \otimes (\cdot): \Bbbk\dcate{Mod} \arrow[shift left, r] & R\dcate{Mod}: \text{lupa} \arrow[shift left, l]
	\end{tikzcd}\]
	ialah endofunktor $L:M\mapsto R\otimes M$ pada $R\dcate{Mod}$. Terdapat
	versi serupa untuk modul-$R$ kanan, yang tidak akan
	diulangi di sini.

	Sekarang terapkan konstruksi bar. Untuk setiap modul-$R$ kiri $M$ kita
	terdapat morfisme kanonik dalam
	$\cate{Ch}_{\geq0}(R\dcate{Mod})$
	% segment-id: o014.li2.chapter8.bar-construction-b.q010
	\[ \epsilon_M: \mathrm{C}(\mathrm{Bar}(M)) \to M; \]
	morfisme ini kuasi-isomorfisme karena
	Korolari~\ref{prop:bar-exactness-G} memastikan bahwa $\epsilon_M$
	kuasi-isomorfisme pada tingkat $\Bbbk\dcate{Mod}$. Secara terperinci,
	% segment-id: o014.li2.chapter8.bar-construction-b.q011
	\[ \mathrm{C}(\mathrm{Bar}(M))_n = L^{n+1}(M) = R^{\otimes (n+1)} \otimes M. \]
	Menurut Contoh~\ref{eg:monad-simplicial} dan deskripsi
	$L\to\identity$ dalam Contoh~\ref{eg:comonad-change-ring} (yakni
	``mengalikan ke dalam''), pemetaan $L^{n+1}(M)\to L^n(M)$ ialah
	% segment-id: o014.li2.chapter8.bar-construction-b.q012
	\begin{multline*}
		\partial_n: r_0 \otimes \cdots \otimes r_n \otimes x \mapsto \\
		\sum_{k=0}^{n-1} (-1)^k \cdots \otimes r_k r_{k+1} \otimes \cdots \otimes x + (-1)^n r_0 \otimes \cdots \otimes r_{n-1} \otimes r_n x,
	\end{multline*}
	dengan $x\in M$ dan $r_i\in R$. Rumus ini tetap bermakna ketika
	$n=0$ ($r_0\otimes x\mapsto r_0x$), dan dengan demikian memberikan
	versi teraugmentasi $\epsilon:\mathrm{C}(\mathrm{Bar}(M))\to M$,
	yaitu $\mathrm{C}^{\mathrm{aug}}(\mathrm{Bar}(M))$.

	Sekarang perhatikan kasus khusus $M=R$. Dalam hal ini
	% segment-id: o014.li2.chapter8.bar-construction-b.q013
	\begin{align*}
		\mathrm{C}(\mathrm{Bar}(R))_n & = R^{\otimes (n+2)} = R \otimes R^{\otimes n} \otimes R, \\
		\partial_n(r_0 \otimes \cdots \otimes r_{n+1}) & = \sum_{k=0}^n (-1)^k \cdots \otimes r_k r_{k+1} \otimes \cdots.
	\end{align*}

	Walaupun $\mathrm{C}(\mathrm{Bar}(R))$ beserta augmentasinya
	$\mathrm{C}^{\mathrm{aug}}(\mathrm{Bar}(R))$ secara definisi merupakan
	kompleks rantai modul-$R$ kiri, deskripsi di atas memperkaya keduanya
	menjadi kompleks rantai bimodul-$(R,R)$: $R$ bekerja melalui
	perkalian kiri pada faktor pertama dan perkalian kanan pada faktor
	terakhir dari $R\otimes R^{\otimes n}\otimes R$. Pengayaan struktur
	ini bukan kebetulan. Karena
	$\mathrm{C}^{\mathrm{aug}}(\mathrm{Bar}(\mathord\cdot))$ adalah funktor,
	setiap endomorfisme $R$ sebagai modul-$R$ kiri (misalnya perkalian
	kanan) terangkat secara alami ke seluruh konstruksi bar.

	Dibandingkan dengan Definisi~\ref{def:bar-Hochschild}, tampak langsung
	bahwa ini persis merupakan kompleks rantai yang digunakan untuk mendefinisikan
	homologi dan kohomologi Hochschild:
	% segment-id: o014.li2.chapter8.bar-construction-b.q014
	\[ \mathrm{C}(\mathrm{Bar}(R)) = \mathsf{B} R, \quad \mathrm{C}^{\mathrm{aug}}(\mathrm{Bar}(R)) = \mathsf{B}' R. \]
	Perhatikan pula bahwa Lema~\ref{prop:HH-bar-exactness} hanyalah kasus
	khusus Teorema~\ref{prop:bar-null-homotopic}.

	Tuliskan $R^e:=R\otimes R^{\opp}$ dan, menurut
	Konvensi~\ref{con:Re-R-R}, identifikasikan modul-$R^e$ kiri,
	bimodul-$(R,R)$, dan modul-$R^e$ kanan. Dari sudut pandang ini, homologi
	Hochschild $\HHm_\bullet(M)$ (atau kohomologi Hochschild
	$\HHm^\bullet(M)$) dari bimodul $M$ merupakan versi turunan dari
	$M\dotimes{R^e}R$ (atau $\Hom_{R^e}(R,M)$): bimodul $R$ diganti oleh
	resolusi barnya $\mathrm{C}(\mathrm{Bar}(R))$, lalu homologi (atau
	kohomologi) dihitung.

	Gagasan penggantian ini muncul di mana-mana dalam studi funktor
	turunan dan kategori turunan, tetapi situasi di sini agak berbeda. Dalam
	kerangka \CHref{sec:cplx} dan
	\CHref{sec:triangulated-derived-cat}, pengganti $R$ seharusnya berupa
	resolusi datarnya (atau resolusi projektifnya) sebagai bimodul. Dengan
	demikian diperoleh $\Tor^{R^e}_\bullet(M,R)$ (atau
	$\Ext^\bullet_{R^e}(R,M)$), ataupun objek yang bersesuaian dalam
	kategori turunan. Contoh~\ref{eg:HH-Tor} (atau
	Contoh~\ref{eg:HH-Ext}) menunjukkan bahwa jika $R$ datar (atau
	projektif) sebagai modul-$\Bbbk$, objek-objek tersebut isomorfik secara
	alami dengan $\HHm_\bullet(M)$ (atau $\HHm^\bullet(M)$); dalam kasus
	umum, hal ini belum tentu benar.
\end{example}

% segment-id: o014.li2.chapter8.bar-construction-b.eg002
\begin{example}[Resolusi bar dan homologi grup]\label{eg:group-H-comonad}
	Misalkan $\Gamma$ sebuah grup. Lanjutkan pembahasan
	Contoh~\ref{eg:HH-comonad} untuk aljabar grup $R:=\Bbbk[\Gamma]$; dalam
	kasus ini $M$ sama dengan sebuah modul-$\Gamma$ kiri dalam arti
	\S\ref{sec:G-mod}. Ingat bahwa
	$(\mathord\cdot)_{\Gamma}:\Gamma\dcate{Mod}\to\Bbbk\dcate{Mod}$ ialah
	funktor koinvarian. Akan ditunjukkan bahwa
	$\epsilon_M:\mathrm{C}(\mathrm{Bar}(M))\to M$ merupakan resolusi
	$(\mathord\cdot)_{\Gamma}$-asiklik dalam $\Gamma\dcate{Mod}$
	(Konvensi~\ref{con:F-acyclic}). Sifat kuasi-isomorfismenya telah
	diketahui. Untuk melihat bahwa setiap
	$\mathrm{Bar}(M)_n=\Bbbk[\Gamma]^{\otimes(n+1)}\otimes M$ bersifat
	$(\mathord\cdot)_{\Gamma}$-asiklik, mula-mula Lema Shapiro
	(Teorema~\ref{prop:Shapiro}) memberikan, untuk setiap modul-$\Bbbk$ $P$,
	% segment-id: o014.li2.chapter8.bar-construction-b.q015
	\[ \Hm_k(\Gamma, \Bbbk[\Gamma] \otimes P ) = \Hm_k\left( \Gamma, \iInd^{\Gamma}_{\{1\}}(P) \right) \simeq \Hm_k(\{1\}, P), \]
	dan ruas kanan selalu $0$ bila $k>0$. Terapkan ini pada
	$P=\Bbbk[\Gamma]^{\otimes n}\otimes M$ untuk memperoleh klaim tersebut.

	Karena itu, Korolari~\ref{prop:acyclic-resolution} mengakibatkan
	$\Hm_n(\Gamma,M)\simeq
	\Hm_n\bigl(\mathrm{C}((\mathrm{Bar}(M)_{\Gamma}))\bigr)$; di sini
	funktor $(\mathord\cdot)_{\Gamma}$ dapat diterapkan baik pada objek
	simplisial maupun, secara ekuivalen, pada kompleks rantai yang
	bersesuaian. Hal ini menafsirkan homologi grup sebagai homologi komonad
	dalam arti latihan bab ini. Lebih tepatnya, mengambil koinvarian-$\Gamma$
	dari $\mathrm{Bar}(M)_n$ sama dengan menghilangkan faktor tensor
	pertama dalam $\Bbbk[\Gamma]^{\otimes(n+1)}\otimes M$ melalui
	$\Bbbk[\Gamma]\twoheadrightarrow\Bbbk$, sehingga diperoleh
	$\Bbbk[\Gamma]^{\otimes n}\otimes M$. Dengan identifikasi ini,
	$\mathrm{Bar}(M)_{n,\Gamma}\to\mathrm{Bar}(M)_{n-1,\Gamma}$ diberikan
	oleh
	% segment-id: o014.li2.chapter8.bar-construction-b.q016
	\begin{multline*}
		g_1 \otimes \cdots \otimes g_n \otimes x \mapsto g_2 \otimes \cdots \otimes g_n \otimes x + \sum_{k=1}^{n-1} (-1)^k \cdots \otimes g_k g_{k+1} \otimes \cdots \otimes x \\
		+ (-1)^n g_1 \otimes \cdots \otimes g_{n-1} \otimes g_n x,
	\end{multline*}
	dengan $g_i\in\Gamma$ dan $x\in M$. Dibandingkan dengan
	Proposisi~\ref{prop:grouph-chain}, tampak langsung bahwa
	$\mathrm{C}(\mathrm{Bar}(M))_{\Gamma}$ persis merupakan kompleks rantai standar
	yang menghitung $\Hm_n(\Gamma,M)$.
\end{example}
